\section{Results}
\label{sec:results}

% SKELETON — fills from generated tables as experiments land.
% Tables in paper/tables/ and macros in paper/generated/numbers.tex are
% produced by scripts/make_paper_tables.py from results/ artifacts.
% NO hand-typed numbers in this section, ever.

\subsection{Detection (M1)}
Detection is essentially solved on this dataset:
\DetMapFifty{} mAP@0.5 (mean of three seeds; std \DetMapFiftyStd)
and \DetMapFiftyNinetyFive{} mAP@0.5:0.95.
Table~\ref{tab:detection} gives per-class AP with supports.

\begin{table}[t]
  \centering
  \caption{Per-class detection AP on the test split (YOLOv8s, seed 0),
  with class supports.}
  \label{tab:detection}
  % AUTO-GENERATED by scripts/make_paper_tables.py — do not edit.
\begin{tabular}{lrrr}
\toprule
class & support & AP@0.5 & AP@0.5:0.95 \\
\midrule
BJT-NPN & 100 & 0.981 & 0.719 \\
BJT-PNP & 73 & 0.965 & 0.714 \\
Capacitor & 315 & 0.995 & 0.660 \\
Diode & 139 & 0.995 & 0.689 \\
GND & 316 & 0.991 & 0.620 \\
I-AC & 69 & 0.991 & 0.814 \\
I-DC & 115 & 0.989 & 0.791 \\
Inductor & 327 & 0.985 & 0.662 \\
MOSFET-N & 54 & 0.889 & 0.670 \\
MOSFET-P & 44 & 0.839 & 0.665 \\
Op-Amp & 116 & 0.994 & 0.795 \\
Resistor & 470 & 0.990 & 0.668 \\
V-AC & 94 & 0.988 & 0.778 \\
V-DC & 154 & 0.994 & 0.805 \\
V-DC (one port) & 102 & 0.995 & 0.641 \\
Wire Crossover & 223 & 0.991 & 0.650 \\
Zener Diode & 75 & 0.985 & 0.693 \\
\bottomrule
\end{tabular}

\end{table}

\subsection{Connectivity: Classical vs.\ Structural (C2 Headline)}
Table~\ref{tab:wire-ablation} reports the progression from the
classical baseline to the full structural configuration. Each step is
validated by a paired per-image bootstrap (Sec.~\ref{sec:metrics}).
\draftnote{narrate: (i) boundary snapping is the single largest jump;
(ii) stitching fixes rail shattering — first non-zero strict success;
(iii) crossover-aware CC was a null result until stitching landed —
stage interactions reverse conclusions (cite the two-run comparison
CSVs); (iv) detection is not the bottleneck.}

\begin{table*}[t]
  \centering
  \caption{The C2 connectivity ablation on the \NumTestImages-image
  verified test split. Bootstrap 95\% CIs in brackets.}
  \label{tab:wire-ablation}
  % AUTO-GENERATED by scripts/make_paper_tables.py — do not edit.
\begin{tabular}{lccccc}
\toprule
configuration & term-pair $F_1$ & net $F_1$ & per-comp.\ acc. & nGED $\downarrow$ & strict \\
\midrule
classical (canny + directional snap) & 0.163 [0.140, 0.190] & 0.433 [0.414, 0.456] & 0.024 [0.011, 0.041] & 0.278 [0.265, 0.291] & 0.000 [0.000, 0.000] \\
+ fixed preprocessing & 0.203 [0.174, 0.231] & 0.431 [0.402, 0.459] & 0.036 [0.018, 0.059] & 0.255 [0.237, 0.272] & 0.000 [0.000, 0.000] \\
+ ink wires + boundary snap & 0.399 [0.363, 0.434] & 0.585 [0.553, 0.617] & 0.142 [0.112, 0.175] & 0.242 [0.223, 0.261] & 0.026 [0.005, 0.053] \\
+ mask-hole stitching & 0.432 [0.395, 0.470] & 0.607 [0.573, 0.639] & 0.174 [0.137, 0.212] & 0.230 [0.211, 0.249] & 0.053 [0.021, 0.084] \\
+ crossover-aware nets (full) & 0.463 [0.425, 0.499] & 0.637 [0.604, 0.666] & 0.184 [0.145, 0.221] & 0.219 [0.200, 0.237] & 0.053 [0.021, 0.084] \\
\bottomrule
\end{tabular}

\end{table*}

\subsection{Full Benchmark, Three Detector Seeds}
Table~\ref{tab:benchmark} reports the primary configuration across
three detector seeds (mean$\pm$std).

\begin{table}[t]
  \centering
  \caption{Primary configuration, three detector training seeds.}
  \label{tab:benchmark}
  % AUTO-GENERATED by scripts/make_paper_tables.py — do not edit.
% 3-seed results not present yet — run scripts/benchmark.py per seed into results/benchmark/seed<N>/.
\multicolumn{2}{c}{\emph{3-seed run pending}}

\end{table}

\subsection{Repair: Solvability Lift with Zero Topology Cost (C5)}
Table~\ref{tab:repair} summarizes the repair layer. DC solvability
rises from \RepSolvBefore{} to \RepSolvAfter{} — a paired per-image
lift with 95\% CI $[\RepLiftCiLo, \RepLiftCiHi]$ and
\RepRegressed{} circuits made worse — at \RepMeanAssumptions{} logged
assumptions per circuit (\RepMeanGauge{} gauge entries). The integrity
rule holds experimentally as well as by construction:
\RepTopoViolations{} of \RepVerifiedImages{} circuits had any component
net assignment altered by repair.

Where the ledger's inferences can be checked against the verified
ground truth, they are. When a GND symbol is present, the reference net
the pipeline selects is the human's ground net in \GroundAccResolved{}
of decidable cases ($n=\GroundNGnd$; \GroundAccStrict{} if the
structurally ambiguous cases are counted as failures rather than
excluded). When no GND symbol is drawn, the ``most connected net''
fallback never matched the human's choice — a small sample, but a clean
negative result that says the fallback is a placeholder, not a method.
\draftnote{add a worked ledger example figure (one circuit, its ledger,
before/after ngspice) — high leverage for the C5 story.}

\begin{table}[t]
  \centering
  \caption{Design-intent repair (C5) on the verified test split.}
  \label{tab:repair}
  % AUTO-GENERATED by scripts/make_paper_tables.py — do not edit.
\begin{tabular}{lc}
\toprule
SPICE syntactic validity & 0.979 \\
DC-solvable before repair & 0.353 \\
DC-solvable after repair & 0.716 \\
mean ledger entries / circuit (assumption) & 4.0 \\
mean ledger entries / circuit (gauge) & 1.7 \\
solvability lift (paired bootstrap 95\% CI) & 0.363 [0.295, 0.432] \\
circuits made worse by repair & 0 \\
topology changed by repair & 0 / 190 \\
ground-choice accuracy, GND symbol present & 0.821 ($n$=183, decidable cases) \\
ground-choice accuracy, no GND symbol & 0.000 ($n$=7) \\
\bottomrule
\end{tabular}

\end{table}

\subsection{Port-Template Localization (C3)}
Table~\ref{tab:ports} reports held-out port-localization error under
two regimes: the true pose (template quality, and the ceiling for any
pose-selection scheme) and axis-only (all a bounding box reveals
without reading the symbol). \draftnote{narrate: pose binning roughly
halves the error for every directional class, and the oracle-vs-axis
gap is precisely what a learned port model would buy; single-port
classes carry no orientation signal at all and are the clearest case
for the learned path.}

\begin{table}[t]
  \centering
  \caption{Held-out port localization from class port templates.
  Error is normalized to the crop width.}
  \label{tab:ports}
  % AUTO-GENERATED by scripts/make_paper_tables.py — do not edit.
\begin{tabular}{lrcccc}
\toprule
& & \multicolumn{2}{c}{oracle pose} & \multicolumn{2}{c}{axis only} \\
\cmidrule(lr){3-4}\cmidrule(lr){5-6}
class & crops & median & $\le$0.10 & median & $\le$0.10 \\
\midrule
BJT-NPN & 825 & 0.134 & 38\% & 0.269 & 23\% \\
BJT-PNP & 633 & 0.128 & 39\% & 0.318 & 17\% \\
Capacitor & 1286 & 0.091 & 56\% & 0.093 & 54\% \\
Diode & 1254 & 0.088 & 58\% & 0.221 & 24\% \\
GND & 1443 & 0.330 & 0\% & 0.330 & 0\% \\
I-AC & 531 & 0.075 & 63\% & 0.078 & 61\% \\
I-DC & 1031 & 0.069 & 65\% & 0.164 & 37\% \\
Inductor & 1391 & 0.069 & 67\% & 0.072 & 65\% \\
MOSFET-N & 386 & 0.169 & 30\% & 0.366 & 10\% \\
MOSFET-P & 386 & 0.198 & 22\% & 0.427 & 3\% \\
Op-Amp & 977 & 0.105 & 47\% & 0.239 & 24\% \\
Resistor & 1459 & 0.069 & 66\% & 0.072 & 63\% \\
V-AC & 642 & 0.066 & 68\% & 0.069 & 67\% \\
V-DC & 1222 & 0.068 & 68\% & 0.191 & 36\% \\
V-DC (one port) & 887 & 0.427 & 0\% & 0.427 & 0\% \\
Wire Crossover & 1359 & 0.109 & 45\% & 0.190 & 14\% \\
Zener Diode & 673 & 0.082 & 60\% & 0.195 & 28\% \\
\bottomrule
\end{tabular}

\end{table}

\subsection{Oracle Stage Attribution (C4)}
The GT-injection oracle attributes end-to-end terminal-pair $F_1$ as
\OracleDetection{} to detection, \OracleWires{} to wire extraction, and
\OracleSnapping{} to terminal snapping, over the
\OracleNValid{} of \NumTestImages{} images whose synthetic wiring
passed verification. \draftnote{narrate carefully — this is a strong
claim and it reverses an earlier internal conclusion: after boundary
snapping landed, snapping is nearly exhausted as an error source and
wire tracing dominates. Also explain the exclusion criterion and why
the waterfall is computed on the valid subset only.}

\subsection{Stratified Performance}
Table~\ref{tab:stratified} breaks results down by drawing
characteristics. \draftnote{narrate: crossing-bearing and
multi-terminal drawings are measurably harder, which is the honest
counterpart to the C2 crossover result — handling crossings helps, and
crossing-heavy drawings are still the hard cases; strict success falls
to zero on the largest circuits because it is a product over
components.}

\begin{table}[t]
  \centering
  \caption{Performance stratified by drawing characteristics.}
  \label{tab:stratified}
  % AUTO-GENERATED by scripts/make_paper_tables.py — do not edit.
\begin{tabular}{lrcccc}
\toprule
stratum & $n$ & term-pair $F_1$ & net $F_1$ & per-comp. & strict \\
\midrule
all & 190 & 0.743 & 0.841 & 0.595 & 0.379 \\
\midrule
small ($\le$8 components) & 72 & 0.875 & 0.918 & 0.786 & 0.681 \\
medium (9-16) & 56 & 0.689 & 0.797 & 0.532 & 0.196 \\
large ($>$16) & 62 & 0.637 & 0.791 & 0.429 & 0.194 \\
has wire crossover & 56 & 0.699 & 0.802 & 0.552 & 0.196 \\
no wire crossover & 134 & 0.761 & 0.858 & 0.613 & 0.455 \\
has 3+ terminal device & 59 & 0.578 & 0.755 & 0.363 & 0.170 \\
only 2-terminal devices & 131 & 0.817 & 0.880 & 0.699 & 0.473 \\
has GND symbol & 183 & 0.740 & 0.839 & 0.590 & 0.377 \\
no GND symbol & 7 & 0.810 & 0.883 & 0.709 & 0.429 \\
\bottomrule
\end{tabular}

\end{table}

\subsection{Failure Cases}
\draftnote{qualitative figure: dense circuits, curved wires, faint
strokes; tie each to the stage the oracle blames. Note when selecting
examples: some apparent total failures were an artifact of GT box
geometry (Sec.~\ref{sec:discussion}) — check against the corrected-box
run before presenting an image as a pipeline failure.}
